\section{System Architecture}\label{sec:system-architecture}

BEACON is designed with a five-layered architecture, with each layer responsible for specific functions.
The layers are as follows:

\subsection{Simulated Physical Layer}\label{subsec:simulated-physical-layer}

The Simulated Physical Layer is responsible for simulating the physical environment and the drones' interactions with it.
The simulated environment comprises a variable bounded mission space, defined by the user, rendered as a water texture with a faint grid overlay.
A central blue circle marks the ``Drone Hub'' in the center of the grid.
Anomalies, such as people in the water, lifeboats, or debris are generated pseudorandomly based on the seed, approximating the distribution of such anomalies in real-world MSAR operations.

\subsection{Individual Drone Autonomy Layer}\label{subsec:individual-drone-autonomy-layer}

Each simulated drone has state variables, such as position, velocity, heading, and battery level.
The drones follow boustrophedon (lawnmower) paths within their assigned Voronoi cells, generated at the Swarm-Level Coordination Layer.
A battery-aware return-to-base system monitors the minimum estimated time to reach the Drone Hub safely.
When the battery level drops below a safe threshold, the drone automatically interrupts its current task and returns to the Drone Hub for recharging.

Each drone also responds to communication degradation, which is simulated as a function of distance from the Drone Hub and the number of drones in the swarm.
When communication degrades, drones will continue to follow their last given path, and will buffer any updates such as newly scanned anomalies.
It will then release this buffered information once communication is re-established, which can be used to update the Digital Twin and the human-swarm interface.

\subsection{Swarm-Level Coordination Layer}\label{subsec:swarm-level-coordination-layer}

The mission-space is partitioned using weighted Voronoi Tessellation, generating cells based on the drones' positions and speeds.
Faster drones receive larger cells, while slower drones receive smaller cells, to ensure that drones complete their assigned search areas at roughly the same time.
Within each cell, boustrophedon sweeps are generated with a custom overlap percentage (15\% by default) to ensure robust coverage.

Cell boundaries act as a first line of collision avoidance, with nominal paths within cells to ensure optimal coverage while avoiding obstacles.

\subsection{Digital Twin \& Decision Support Layer}\label{subsec:digital-twin-decision-support-layer}

A central state management system maintains the state of the simulated drones, through drone states, anomaly states, paths, and mission parameters.
This layer computes derived indicators such as alert burden, comms resilience, and coverage efficiency, which are then used to generate alerts and notifications for the human-swarm interface.
While this layer is currently simulated, it is intended to be connected to real-time data from physical drones in future interactions, to create a full Digital Twin system.

\subsection{Human-Swarm Interface Layer}\label{subsec:human-swarm-interface-layer}

The top-level interface comprises a central top-down map view of the mission space, showing drones, anomalies (if not hidden), coverage paths and assigned cells if generated.
Above the map, there is also a metrics dashboard showing key indicators such as coverage efficiency, alert burden, and comms resilience.
Under the map, alerts and notifications are shown in a feed, following a model similar to the flight deck alerting systems in Crew Alerting Systems used in aircraft.
This means that alerts are prioritized based on severity, such as Advisories, Amber Cautions, Red Warnings, and Status Alerts, similar to the Central Alerting Systems used in Aviation~\cite{RN3}.

These alerts are generated based on the state of the drones, such as low battery, communication degradation, or anomaly detection.
The alerts are then accompanied by an associated chime, and a visual notification on the map, to ensure that operators are aware of critical events.

Prior to mission start, the user can manage scenarios, pick through a list of pre-generated scenarios, or generate a custom scenario by inputting parameters.
After the mission, they can also review the mission's performance, through a debriefing screen showing key metrics, and an optional NASA-TLX workload assessment form.

The layout and the appearance are informed by HCI principles, which will be further detailed in the following section.
